\documentclass[10pt,a4paper]{amsart}
\usepackage[margin=19mm]{geometry}
\usepackage{amsmath,amssymb,mathtools}
\usepackage[scr=rsfs,cal=euler]{mathalfa}
\usepackage{xfrac}
\usepackage[unicode,hidelinks,bookmarks=false]{hyperref}
\newcommand{\ie}{i.e.\ }
\newcommand{\cf}{cf.\ }
\newcommand{\resp}{respectively\ }
\newcommand{\Subsection}{\subsection}
\providecommand{\nobreakdash}{\nobreak\hbox{-}}
\setlength{\emergencystretch}{2em}
\multlinegap0pt
\usepackage{bm}
\usepackage{bbm}
\usepackage{dsfont}
\usepackage{xspace}
\usepackage{cancel}
\usepackage{mathtools}
\usepackage{cleveref}
\usepackage{amsthm}
\makeatletter
\@ifundefined{assumptions}{\newtheorem{assumptions}{Assumptions}}{}
\@ifundefined{theorem}{\newtheorem{theorem}{Theorem}}{}
\@ifundefined{claim}{\newtheorem{claim}[theorem]{Claim}}{}
\@ifundefined{definition}{\newtheorem{definition}[theorem]{Definition}}{}
\makeatother
\newcommand{\R}{\mathbb{R}}
\newcommand{\bfb}{\boldsymbol{b}}
\newcommand{\bfh}{\boldsymbol{h}}
\newcommand{\bftheta}{\boldsymbol{\theta}}
\newcommand{\bfu}{\boldsymbol{u}}
\newcommand{\bfv}{\boldsymbol{v}}
\newcommand{\bfw}{\boldsymbol{w}}
\newcommand{\calA}{\mathcal{A}}
\newcommand{\norm}[1]{\left\lVert#1\right\rVert}
\title[Rigorous Asymptotics for First-Order Algorithms Through the ]{Rigorous Asymptotics for First-Order Algorithms Through the Dynamical Cavity Method}
\author{Yatin Dandi, David Gamarnik, Francisco Pernice, Lenka Zdeborov\'a}
\date{Source: arXiv:2603.14573v1 (CC BY 4.0)}
\begin{document}
\maketitle

\noindent\emph{Source and licence.} Rigorous Asymptotics for First-Order Algorithms Through the Dynamical Cavity Method. Version arXiv:2603.14573v1 (CC BY 4.0), distributed under the Creative Commons Attribution 4.0 International licence, as recorded in the arXiv licence metadata for this exact version and shown on the abstract page https://arxiv.org/abs/2603.14573v1. Full rights notice: licenses/arxiv-2603.14573-CC-BY-4.0.txt.\par\smallskip

\noindent\textit{Editorial scope.} Counted unit: the Claim whose source label is \texttt{claim:recursive-claim-1} in the article (source0.tex lines 957--963, source label \texttt{claim:recursive-claim-1}), delivered together with the article's own complete original proof of it (source0.tex lines 964--1059, one \texttt{proof} environment of 96 lines). No mathematics is written, repaired or re-derived: statement and proof are the article's own text line for line. Required context retained uncounted: defn:GFOM (lines 219--230); assump:2 (lines 284--291); assump:4 (lines 284--291); assump:1-weak (lines 483--490); thm:derivative-bounds (lines 493--509); eq:V-FTC (lines 654--656). Every definition, display and label of those lines is retained unchanged. Omitted from this package and not redistributed: 1--218, 231--283, 292--482, 491--492, 510--653, 657--956, 1060--1290. Nothing inside a retained excerpt is omitted or altered: every retained body is delivered exactly as the article's lines are written, so no omission receipt of any kind accompanies this package. Citations: the retained excerpts carry no citation command at all, so no bibliography entry of the article is transcribed and no reference is redistributed. Reference adaptation: none was needed and none was made; every reference inside the delivered unit resolves inside it, so no unresolved reference or citation is produced. Printed slips kept literal: no printed word, symbol, hypothesis or number of the counted statement or of its proof was altered. Layout and class: the article's own class is \texttt{article} with packages including graphicx, amsmath, amssymb, mathtools, algorithm, algpseudocode, geometry, color, amsthm, biblatex, tikz, hyperref, cleveref, enumitem; the wrapper is the lane's pinned \texttt{amsart} editorial wrapper at 10pt on A4, which restates the theorem-like environments the retained text needs and adds the article's own macro definitions that the retained text needs (\texttt{\textbackslash R}, \texttt{\textbackslash bfb}, \texttt{\textbackslash bfh}, \texttt{\textbackslash bftheta}, \texttt{\textbackslash bfu}, \texttt{\textbackslash bfv}, \texttt{\textbackslash bfw}, \texttt{\textbackslash calA}, \texttt{\textbackslash norm}), with extra wrapper packages (bm, bbm, dsfont, xspace, cancel, mathtools, cleveref). The delivered numbering is the unit's own. Assets: no retained span contains a figure environment, a graphics-inclusion command, a PDF-page inclusion, a table or a TikZ picture; no image file is copied into this package, the package declares no asset dependency and no image file is redistributed with it. Licence: version arXiv:2603.14573v1 of this article is distributed under the Creative Commons Attribution 4.0 International licence, as recorded in the arXiv licence metadata for this exact version and shown on the abstract page https://arxiv.org/abs/2603.14573v1. A full rights notice is delivered at licenses/arxiv-2603.14573-CC-BY-4.0.txt. Characters: every retained excerpt is pure ASCII, so the delivered engine, XeLaTeX with its default Latin Modern text font, reports no missing character.
\begin{definition}\label{defn:GFOM}
    Let $X \in \R^{N\times d}, \bftheta, \bfv^0 \in \R^d,\bfb^* = X\bftheta \in \R^N.$ For $t=1,\dots,$ suppose $\bfw^t \in \R^N, \tilde{\bfw}^t \in \R^d$, and $F_t, G_t:\R^{2t + 1}\to \R$. The \emph{GFOM orbit} is the iteration
    \begin{align}\label{eq:GFOM-orbit}
        \begin{cases}
            \bfb^t = X \bfv^{t-1} &\in \R^N \\
            \bfu^t = F_t(\bfb^*, \bfb^1,\dots,\bfb^t; \bfw^1,\dots,\bfw^t) & \in \R^N \\
            \bfh^t = X^T \bfu^t &\in \R^d \\
            \bfv^t = G_t(\bfh^1,\dots,\bfh^t;\bfv^0,\tilde{\bfw}^1, \dots, \tilde{\bfw}^t) & \in \R^d.
        \end{cases}
    \end{align}
    In our notation, we will often write $\bfu^t = F_t(\bfb^*, \bfb^1,\dots,\bfb^t)$ and $\bfv^t = G_t(\bfh^1,\dots,\bfh^t),$ making the dependence on $\bfw,\tilde{\bfw}$ and $\bfv^0$ implicit.
\end{definition}

\begin{assumptions}
  \item\label{assump:1}
  The entries of $X$ are sampled i.i.d.\ from a $\sigma/\sqrt{N}$-subgaussian measure with mean zero and
  variance $1/N$.
  \item\label{assump:2} The entries of $\bfv^0,\bftheta \in \R^d$ are i.i.d. $(\theta_j,v^0_j)\sim\nu_0$ and independent of $X$, where $\nu_0$ is $\sigma$-subgaussian.
  \item \label{assump:3} The functions $F_t,G_t$ are $L$-Lipschitz, where $L$ is a constant independent of $N.$
  \item \label{assump:4} The noise vectors $\bfw^t,\tilde{\bfw}^t$ are independent of each other and of everything else, with entries i.i.d. $N(0,1).$
\end{assumptions}

\begin{assumptions}
  \item[(A1')]\makeatletter
    \protected@edef\@currentlabel{A1'}%
  \makeatother
  \label{assump:1-weak}
  The entries of $X$ are sampled i.i.d.\ from a $\sigma/\sqrt{N}$-subgaussian
  measure with mean zero.
\end{assumptions}

\begin{theorem}\label{thm:derivative-bounds}
    Suppose assumptions \eqref{assump:1-weak} and \eqref{assump:2}-\eqref{assump:4} hold in \Cref{defn:GFOM}, and $F_t,G_t$ are $C^\infty_{1b}.$  Then, for every $t,p\geq 1$ and $m \geq 0$, there exists a constant $\Gamma_{t,m,p}$ independent of $N$ such that the following hold for every multiset of derivatives $P$ such that $|P|=m$ and for every $i \in [N]$ and $j \in [d]$.
    \begin{enumerate}
        \item For $s=1,\dots,t,$ we have
        \begin{align*}
            \max\{\norm{\partial_P b^*_i}_p,\; \norm{\partial_P b^s_i}_p, \; \norm{\partial_P u^s_i}_p,\; \norm{\partial_P h^s_j}_p,\; \norm{\partial_P v^s_j}_p \} &\leq \Gamma_{t,m,p}.
        \end{align*}
        \item If $i$ is not all reaching in $P,$ then
        \begin{align*}
            \max\{\norm{\partial_P b^*_i}_p,\; \norm{\partial_P b^s_i}_p, \; \norm{\partial_P u^s_i}_p\} &\leq \Gamma_{t,m,p}/\sqrt{N}.
        \end{align*}
        Similarly, if $j$ is not all reaching in $P$, then
        \begin{align*}
            \max\{\norm{\partial_P h^s_j}_p,\; \norm{\partial_P v^s_j}_p \} &\leq \Gamma_{t,m,p}/\sqrt{N}.
        \end{align*}
    \end{enumerate}
\end{theorem}

\begin{align}
    V_{j_1,\dots,j_p}(X) &= V'_{j_1,\dots,j_p}(X) + \int_{[0,1]^q}  \partial_{X_{ij_1'}}\dots \partial_{X_{ij_p'}}V_{j_1,\dots,j_p}(X^{i,\{j'_1,\dots,j'_q\}}(\eta_1,\dots,\eta_q))\prod_{\ell=1}^q (X_{i,j'_\ell}d\eta_\ell), \label{eq:V-FTC}
\end{align}

\begin{claim}\label{claim:recursive-claim-1}
    For every $t,p\geq 1,i\in [N]$ and $j\in [d],$ we have 
    \begin{align*}
        \lim \norm{b_i^t - \tilde{b}_i^t} &= 0 \\
        \lim \norm{h_j^t - \tilde{h}_j^t} &= 0.
    \end{align*}
\end{claim}

\begin{proof}[of \Cref{claim:recursive-claim-1}]
      The proof is again by induction. At each step of the induction, we assume that for all $s,t\leq T-1$, all $p\geq 1$ and all $i\in [N],j\in [d],$ we have proved that
    \begin{align*}
        \max\left\{\norm{ h_j^t - \tilde{ h}_j^t}_p , \norm{\partial_{ h_j^s} h_j^t - \tilde{\partial}_{\tilde{ h}_j^s}\tilde{ h}_j^t}_p\right\}&\lesssim 1/\sqrt{N} \\
        \max\left\{\norm{ v_j^t - \tilde{ v}_j^t}_p , \norm{\partial_{ h_j^s} v_j^t - \tilde{\partial}_{\tilde{ h}_j^s}\tilde{ v}_j^t}_p\right\}&\lesssim 1/\sqrt{N} \\
        \max\left\{\norm{ b_i^t - \tilde{ b}_i^t}_p , \norm{\partial_{ b_i^s} b_i^t - \tilde{\partial}_{\tilde{ b}_i^s}\tilde{ b}_i^t}_p,\norm{\partial_{ b^*_i} b_i^t - \tilde{\partial}_{\tilde{ b}^*_i}\tilde{ b}_i^t}_p \right\}&\lesssim 1/\sqrt{N} \\
        \max\left\{\norm{ u_i^t - \tilde{ u}_i^t}_p , \norm{\partial_{ b_i^s} u_i^t - \tilde{\partial}_{\tilde{ b}_i^s}\tilde{ u}_i^t}_p,\norm{\partial_{ b_i^*} u_i^t - \tilde{\partial}_{\tilde{ b}^*_i}\tilde{ u}_i^t}_p \right\}&\lesssim 1/\sqrt{N}
    \end{align*}
    and aim to prove the same for $s<t=T$ (since in the case $s=t$, all partial derivatves are equal to 1 so the statement is trivial). Note that the base case $T=0$ is trivially true. Since all inductive cases are analogous, we elaborate only the cases of $\bfh$ and $\bfv.$ For $\bfh$, we have
    \begin{align*}
        \norm{h_j^t - \tilde{h}_j^t}_p &= \norm{\sum_{i=1}^N X_{ij}^2 \left(\tilde{\theta}_j \tilde{\partial}_{\tilde{b}_i^*} \tilde{u}_i^t - \theta_j \partial_{b_i^*}u_{i/ij}^t + \sum_{s=1}^t (\tilde{v}_j^{s-1} \tilde{\partial}_{\tilde{b}_i^s} \tilde{u}_i^t - v_{j/ij}^{s-1} \partial_{b_i^s}u_{i/ij}^t )\right)}_p + \norm{\Delta_j^t}_p \\
        &\lesssim  \frac{1}{N}\sum_{i=1}^N\left(\norm{ \tilde{\theta}_j \tilde{\partial}_{\tilde{b}_i^*} \tilde{u}_i^t - \theta_j \partial_{b_i^*}u_{i/ij}^t }_{2p}+ \sum_{s=1}^t \norm{\tilde{v}_j^{s-1} \tilde{\partial}_{\tilde{b}_i^s} \tilde{u}_i^t - v_{j/ij}^{s-1} \partial_{b_i^s}u_{i/ij}^t )}_{2p}\right) + \norm{\Delta_j^t}_p.
    \end{align*}
    We can bound each term individually. For $\norm{\Delta_j^t}_p$, using \Cref{thm:derivative-bounds}, we have 
    \begin{align*}
        \norm{\Delta_j^t}_p &\lesssim  \sum_{s=1}^{t-1} \sum_{i=1}^N \norm{u^{s}_{i/ij}X_{ij}^2\partial_{h_j^s}u_{i/ij}^{t}   }_p+ \frac{1}{2} \int_0^1 d\eta(1-\eta) \sum_{i=1}^N \norm{X_{ij}^3 \partial^2_{X_{ij}}u_i^{t}(X^{ij}(\eta))}_p \\
        &\lesssim \sum_{s=1}^{t-1} \sum_{i=1}^N \frac{1}{N}\norm{\partial_{h_j^s}u_{i/ij}^{t}   }_{3p}+ \frac{1}{2} \int_0^1 d\eta(1-\eta) \sum_{i=1}^N \frac{1}{N^{3/2}}\norm{\partial^2_{X_{ij}}u_i^{t}(X^{ij}(\eta))}_{2p} \\
         &\lesssim \sum_{s=1}^{t-1} \sum_{i=1}^N \frac{1}{N^{3/2}}+ \frac{1}{2} \int_0^1 d\eta(1-\eta) \sum_{i=1}^N \frac{1}{N^{3/2}} \\
         &\lesssim 1/\sqrt{N}.
    \end{align*}
    Similarly, for $\norm{ \tilde{\theta}_j \tilde{\partial}_{\tilde{b}_i^*} \tilde{u}_i^t - \theta_j \partial_{b_i^*}u_{i/ij}^t }_{2p}$ we have
    \begin{align*}
        \norm{ \tilde{\theta}_j \tilde{\partial}_{\tilde{b}_i^*} \tilde{u}_i^t - \theta_j \partial_{b_i^*}u_{i/ij}^t }_{2p} &= \norm{ \theta_j \tilde{\partial}_{\tilde{b}_i^*} \tilde{u}_i^t - \theta_j \partial_{b_i^*}u_{i/ij}^t }_{2p} \\
        &\lesssim \norm{  \tilde{\partial}_{\tilde{b}_i^*} \tilde{u}_i^t - \partial_{b_i^*}u_{i/ij}^t }_{4p} \\
        &\leq \norm{  \tilde{\partial}_{\tilde{b}_i^*} \tilde{u}_i^t - \partial_{b_i^*}u_{i}^t }_{4p} + \norm{\partial_{b_i^*}u_{i}^t - \partial_{b_i^*}u_{i/ij}^t}_{4p}.
    \end{align*}
    The first term is $O( 1/\sqrt{N})$ by inductive hypothesis. For the second term, we can write 
    \begin{align}
        \norm{\partial_{b_i^*}u_{i}^t - \partial_{b_i^*}u_{i/ij}^t}_{4p} &\leq \int_0^1d\eta \norm{X_{ij} \partial_{X_{ij}}\partial_{b_i^*}u_{i}^t(X^{ij}(\eta)}_{4p}\label{eq:taylor-expansion-to-put-xij-back} \\
        &\lesssim 1/\sqrt{N}\nonumber
    \end{align}
    again by \Cref{thm:derivative-bounds}. A similar proof to the above shows that 
    \begin{align*}
        \norm{\tilde{v}_j^{s-1} \tilde{\partial}_{\tilde{b}_i^s} \tilde{u}_i^t - v_{j/ij}^{s-1} \partial_{b_i^s}u_{i/ij}^t )}_{2p} &\lesssim 1/\sqrt{N}
    \end{align*}
    as well.

    Next we show that we have 
    \begin{align*}
        \norm{\partial_{ h_j^s} h_j^t - \tilde{\partial}_{\tilde{ h}_j^s}\tilde{ h}_j^t}_p &\lesssim 1/\sqrt{N}.
    \end{align*}
    Here we can write 
    \begin{align}
        &\norm{\partial_{ h_j^s} h_j^t - \tilde{\partial}_{\tilde{ h}_j^s}\tilde{ h}_j^t}_p \nonumber\\
        &\leq \norm{\partial_{ h_j^s} g_j^t}_p + \norm{\sum_{i=1}^N X_{ij}^2\left(\tilde{\partial}_{\tilde{ h}_j^s}\tilde{\theta}_j \tilde{\partial}_{\tilde{b}_i^*} \tilde{u}_i^t - \partial_{ h_j^s}\theta_j \partial_{b_i^*}u_{i/ij}^t + \sum_{s=1}^t (\tilde{\partial}_{\tilde{ h}_j^s}\tilde{v}_j^{s-1} \tilde{\partial}_{\tilde{b}_i^s} \tilde{u}_i^t - \partial_{ h_j^s}v_{j/ij}^{s-1} \partial_{b_i^s}u_{i/ij}^t )\right)}_p \nonumber\\
        &\qquad \qquad + \norm{\sum_{i=1}^N X_{ij}^2\left( \theta_j \partial_{ h_j^s}\partial_{b_i^*}u_{i/ij}^t + \sum_{s=1}^t v_{j/ij}^{s-1} \partial_{ h_j^s}\partial_{b_i^s}u_{i/ij}^t \right)}_p + \norm{\partial_{ h_j^s} \Delta_j^t}_p \nonumber\\
        &\lesssim \norm{\partial_{ h_j^s} g_j^t}_p + \frac{1}{N}\sum_{i=1}^N \left(\norm{\tilde{\partial}_{\tilde{ h}_j^s}\tilde{\theta}_j \tilde{\partial}_{\tilde{b}_i^*} \tilde{u}_i^t - \partial_{ h_j^s}\theta_j \partial_{b_i^*}u_{i/ij}^t}_{2p} + \sum_{s=1}^t \norm{\tilde{\partial}_{\tilde{ h}_j^s}\tilde{v}_j^{s-1} \tilde{\partial}_{\tilde{b}_i^s} \tilde{u}_i^t - \partial_{ h_j^s}v_{j/ij}^{s-1} \partial_{b_i^s}u_{i/ij}^t }_{2p} \right) \nonumber\\
        &\qquad \qquad + \frac{1}{N}\sum_{i=1}^N\left(\norm{ \theta_j \partial_{ h_j^s}\partial_{b_i^*}u_{i/ij}^t}_{2p} + \sum_{s=1}^t\norm{ v_{j/ij}^{s-1} \partial_{ h_j^s}\partial_{b_i^s}u_{i/ij}^t }_{2p} \right)+ \norm{\partial_{ h_j^s} \Delta_j^t}_p \label{eq:h-partials-difference} 
    \end{align}
    The second and last terms above can be bounded in a completely analogous way to how we bounded the corresponding terms in $\norm{h_j^t - \tilde{h}_j^t}_p$. The summands in the third term above can similarly be bounded since e.g.
    \begin{align*}
        \norm{ \theta_j \partial_{ h_j^s}\partial_{b_i^*}u_{i/ij}^t}_{2p} &\lesssim \norm{  \partial_{ h_j^s}\partial_{b_i^*}u_{i/ij}^t}_{4p}\\
        &\lesssim 1/\sqrt{N}
    \end{align*}
    by \Cref{thm:derivative-bounds}. What remains is to bound the first term $ \norm{\partial_{ h_j^s} g_j^t}_p.$ To bound this term, we use the same Taylor expansion as in \eqref{eq:V-FTC}. In this case we set 
    \begin{align}\label{eq:V-expansion-new}
    V_{i_1,\dots,i_p}(X) &:= (\partial_{ h_j^s} u_{i_1/i_1j}^{t-1}) \cdots (\partial_{ h_j^s} u_{i_p/i_pj}^{t-1}).
\end{align}
Following the same argument as in the proof of \Cref{thm:derivative-bounds}, we obtain 
\begin{align}
    &\norm{\partial_{h_j^s} g_j^t}_p^p \nonumber\\
    &\lesssim  \sum_{q=0}^p \int_{[0,1]^q} \prod_{\ell=1}^q d\eta_\ell\frac{1}{N^{(p+q)/2}}\sum_{(i_1,\dots,i_p) \in \calA_q} \norm{\partial_{X_{i'_1j}}\dots \partial_{X_{i'_qj}} V_{i_1,\dots,i_p}(X^{\{i'_1,\dots,i'_q\},j}(\eta_1,\dots,\eta_q))}_{p+1}. \nonumber
\end{align}
We claim that, for each of the terms in the sum above, we have 
\[
\norm{\partial_{X_{i'_1j}}\dots \partial_{X_{i'_qj}} V_{i_1,\dots,i_p}(X^{\{i'_1,\dots,i'_q\},j}(\eta_1,\dots,\eta_q))}_{p+1} \lesssim (1/\sqrt{N})^p.
\]
To prove this, it suffices to show that, for any way to distribute the partial derivatives $\partial_{X_{i'_1j}},\dots, \partial_{X_{i'_qj}}$ among the factors in $(\partial_{ h_j^s} u_{i_1/i_1j}^{t-1}), \cdots ,(\partial_{ h_j^s} u_{i_p/i_pj}^{t-1})$, every factor must have moments that are of order  $1/\sqrt{N}$. Note first that, since $X_{i'_k,j}$ has been set to zero in $\partial_{h_j^s}u_{i_k'/i_k'j}^{t-1},$ we can assume that no derivative $\partial_{X_{i'_k,j}}$ gets distributed to $\partial_{h_j^s}u_{i_k'/i_k'j}^{t-1},$ since otherwise the corresponding term will be equal to zero. But recall that the $i'_1,\dots,i'_k$ are unique among $i_1,\dots,i_p.$ Hence, any factors that get any derivatives at all distributed to them must have moments that are $\lesssim 1/\sqrt{N}$ by \Cref{thm:derivative-bounds}. Morever, any term that doesn't get any derivatives distributed to it is still being acted on by $\partial_{h_j^s}$, and hence again has moments that are $O(1/\sqrt{N})$ by \Cref{thm:derivative-bounds}. This proves that
\begin{align*}
    \norm{\partial_{h_j^s} g_j^t}_p^p &\lesssim \sum_{q=0}^p \frac{1}{N^{(p+q)/2}} \cdot |\calA_q|\cdot \frac{1}{N^{p/2}} \\
    &\lesssim \frac{1}{N^{p/2}}.
\end{align*}
Plugging this back into \eqref{eq:h-partials-difference}, we obtain 
\begin{align*}
    \norm{ \partial_{h_j^s}h_j^t - \tilde{\partial}_{\tilde{h}_j^s}\tilde{h}_j^t }_p &\lesssim 1/\sqrt{N},
\end{align*}
as desired.

To conclude the proof of \Cref{claim:recursive-claim-1}, we elaborate the inductive case for $\bfv.$ By Lipschitzness of $G_t$, we have 
\begin{align*}
    \norm{v_j^t - \tilde{v}_j^t}_p &\lesssim \sum_{s=1}^{t}\norm{h_j^s - \tilde{h}_j^s}_p \\
    &\lesssim 1/\sqrt{N}
\end{align*}
by inductive hypothesis. Similarly, by chain rule we have we have 
\begin{align*}
&\norm{\partial_{h_j^s} v_j^t - \tilde{\partial}_{\tilde{h}_j^s} \tilde{v}_j^t}_p \\
&= \sum_{k=1}^t \norm{\partial_k G_t(h_j^1, \dots, h_j^t; v_j^0,\tilde{w}_j^1,\dots, \tilde{w}_j^t) \partial_{h_j^s} h_j^k - \partial_k G_t(\tilde{h}_j^1, \dots, \tilde{h}_j^t; v_j^0,\tilde{w}_j^1,\dots, \tilde{w}_j^t) \tilde{\partial}_{\tilde{h}_j^s} \tilde{h}_j^k  }_p \\
&\leq \sum_{k=1}^t \norm{\partial_k G_t(h_j^1, \dots, h_j^t; v_j^0,\tilde{w}_j^1,\dots, \tilde{w}_j^t) \partial_{h_j^s} h_j^k - \partial_k G_t(\tilde{h}_j^1, \dots, \tilde{h}_j^t; v_j^0,\tilde{w}_j^1,\dots, \tilde{w}_j^t) \partial_{h_j^s} h_j^k}_p \\
&+ \sum_{k=1}^t \norm{\partial_k G_t(\tilde{h}_j^1, \dots, \tilde{h}_j^t; v_j^0,\tilde{w}_j^1,\dots, \tilde{w}_j^t) \partial_{h_j^s} h_j^k- \partial_k G_t(\tilde{h}_j^1, \dots, \tilde{h}_j^t; v_j^0,\tilde{w}_j^1,\dots, \tilde{w}_j^t) \tilde{\partial}_{\tilde{h}_j^s} \tilde{h}_j^k  }_p  \\
&\lesssim \sum_{k=1}^t \norm{\partial_k G_t(h_j^1, \dots, h_j^t; v_j^0,\tilde{w}_j^1,\dots, \tilde{w}_j^t) - \partial_k G_t(\tilde{h}_j^1, \dots, \tilde{h}_j^t; v_j^0,\tilde{w}_j^1,\dots, \tilde{w}_j^t) }_{2p}\\
&+ \sum_{k=1}^t \norm{ \partial_{h_j^s} h_j^k-  \tilde{\partial}_{\tilde{h}_j^s} \tilde{h}_j^k  }_{2p} \\
&\lesssim \sum_{k=1}^t\sum_{k'=1}^t \norm{h_j^{k'} - \tilde{h}_j^{k'}}_{2p} + \sum_{k=1}^t \norm{ \partial_{h_j^s} h_j^k-  \tilde{\partial}_{\tilde{h}_j^s} \tilde{h}_j^k  }_{2p} \\ 
&\lesssim 1/\sqrt{N}
\end{align*} 
by \Cref{thm:derivative-bounds} and the inductive hypothesis. This concludes the proof of \Cref{claim:recursive-claim-1}.
\end{proof}

\end{document}
